\section{Introduction}
In article \cite{SalaMartinDoppelhoferMiller_BACE} authors suggest a
technique for averaging the parameter estimates -- Bayesian Averaging
of Classical Estimates (BACE), which enables to measure the importance of
particular potential regressors. This approach is alternative to the
familiar and earlier-applied Bayesian Model Averaging (BMA) technique,
from which it differs mainly because of its less restrictive \textit{prior}
assumptions. The authors showed that BACE weighting method may
be understood as a limiting case of BMA approach  when the \textit{prior}
information is dominated by the data. The averaging of the parameter
estimates is performed across all possible combinations of linear regression models
obtained by means of Ordinary Least Squares (OLS)
 with classical assumptions.

The BACE approach has several interesting features in relation to
the techniques such as BMA or robustness-checking methods (\textcolor{red}{reference}).
First of all, it is not necessary to make assumptions concerning
\textit{prior} distribution for parameters but only
one simple assumption, on the expected model size.

An example of the application of the BACE technique for forecasting
macroeconomic indicators can be found, among others, in the articles
\cite{vanDijk2004Forecasting, AlbisMapa2014MPRA55902, BialowolskiKuszewskiWitkowski2015Empirica}.